\chapter{Oscillator-Processor Duality}
\label{ch:processor}

\papernote{Source paper: \emph{Categorical Processing Unit} --- \texttt{mekaneck/docs/processor/}}

\begin{chapterabstract}
The oscillator-processor duality theorem states that any oscillator with angular frequency $\omega$ functions as a computational processor with processing rate $R_{\mathrm{compute}} = \omega/(2\pi)$. This is not a metaphor but a mathematical identity: the oscillatory cycle \emph{is} the computational cycle, and categorical state transitions during the cycle \emph{are} computations. The entropy reformulation $S = f(\omega_{\mathrm{final}}, \phi_{\mathrm{final}}, A_{\mathrm{final}})$ enables $O(1)$ complexity navigation of state space. Biological semiconductor implementation: oscillatory holes function as p-type carriers; universal quantum gates are realised at 758~Hz with 10~ms coherence times.
\end{chapterabstract}

\input{../mekaneck/docs/processor/categorical-processing-unit.tex}
